\section{修正后的判断}

导言里悬置的问题，到这里可以给出修正后的判断了。所有判断按本文的置信度分级标注，并附与第 2 至 10 章的对应依据。

\textbf{判断一：技术层面，若手稿成立，这是一次问题层面的范式级推进。}\cfH{}三窄度区间 + 地平线的组合，是对四维特有障碍（低次代数曲面上的聚集）的第一份系统性回应；它放弃了门内“形态穷举”的路线，改用参数分层与极值改进，方法的完整度从第 5 章的骨架重建来看是自洽的。这里的“范式级”限定在问题内部——它替换的是解决四维挂谷的组织方式，而不是数学思想的坐标系；三层分解见判断七。但“若成立”三个字承载的分量很大——本判断的可信度取决于验证状态，而不是文本的精致程度。

\textbf{判断二：谱系层面，这是挂谷门派有史以来距离最远的一次跨门派借入。}\cfH{}熵账本、概率耦合、实代数几何纪律与乘法卷积 Fourier 衰减在门内文献中零命中（第 7.3 节词频表），这一事实不依赖于对手稿正确性的判断。历史经验（加法组合、多项式方法的消化各用了一代人）提示：借入距离越远，理解成本越高，而理解成本正是当下最大的隐性支出。

\textbf{判断三：对数学共同体，验证瓶颈已取代解题难度成为主要矛盾。}\cfH{}一个可以从集群侧无限提速的生成端，对接一个只能线性扩展的验证端，“生成快、验证慢”是结构性不对称；074 族是这个不对称的第一次标志性显影——272 页、零形式化、零审评，而全世界的合格审读者只有两位数。

\textbf{判断四：形式化是部分答案，且只覆盖机械层。}\cfM{}它在依赖链核验上不可替代，人类侧的三维挂谷形式化已经证明了可行性；但它既不能担保自然语言层与形式层的语义一致，也不能产生概念理解。把它当作验证的全部，或当作可以忽略的余事，都会误判成本结构。

\textbf{判断五：对读者与青年学者的使用建议。}\cfM{}把这两篇稿件当作一份“待核验的研究计划书 + 工程样机”，而不是定理库中的新条目：先读第 5 章的骨架、附录 B 的清单，再决定投入方向；核验的副产品（E1--E5 类小课题）本身具有独立发表价值。任何“相信”或“不信”的整体表态，都不如先完成链条根部的一次复演。

\textbf{判断六：对“AI 辅助数学”这一宏观问题的姿态。}\cfM{}本案例既支持乐观叙事（方法论创新密度确实高于同期人类论文的平均值），也支持悲观叙事（验证与理解的欠账被同步放大）。两种叙事共享一个事实基础：数学共同体的核心资产不是结论库，而是“理解—复用—再生成”的循环；批量产出对这个循环的冲击，远大于对结论库的冲击——这也正是 AHM 声明与 OpenAI 发布口径之间争议的真正落点\upcite{ahm2026,openairelease2026,tao2026}。

\textbf{判断七（“开创还是闭包”的分层定位）。}\cfM{}把“是否发明了新数学思想”拆成三层来问，答案才站得住。\emph{命题层}：$\dim_{\mathrm{H}}K=4$ 若成立，属于发现——它陈述的是对象本来具有的性质，无所谓发明，只有真伪。\emph{装置层}：图卡、地平线、比较原理、次数恢复、双帧实验、保边回填这些对象与机制属于发明——它们在门内文献中零命中（第 7.3 节词频证据），但每一项都有可辨认的亲缘（多项式 Remez、划分、极值反证、熵、耦合），是“在既有语言内新造机件”。\emph{范式层}：极值反证加局部拟合加逐尺度归纳的大框架延续自门内传统；本次真正带思想性质的新意，是把“局部理解所需付出的代价”形式化为可极值化的资源（地平线），并首次把信息论、分形几何与实代数几何三套语言组装进同一套证明纪律。据此，本文的定位是：\textbf{它不是“现有知识方法的凸闭包推论”}——该闭包在 $3.059$ 与 $4$ 之间停摆了二十余年，且证明需要新增对象；\textbf{也够不上“革命性的数学思想创新”}——它没有像多项式方法或熵方法当年那样重写整个领域的问题意识。最准确的刻度是：\emph{继承框架、发明装置、跨门派综合的破局性工作}——对问题本身是决定性跃迁，对数学思想是高水平综合，对 AI 参与数学的工具箱则是第一次“全域架构级”产出。以上以手稿成立为前提；若最终被修正，装置层的评估需按存留的有效部件重新计算。三层结论的技术底稿分别为：装置的精确定义与推导细节（附录 C.1）、与传统极值反证的形式化对比（附录 C.2，含 $H(E)$ 与改进引理的公式化写法）、三套外门派语言的映射表（附录 C.3）。

回到第 2 章的问题本身：方向完备性与空间压缩的张力，在四维以上还没有被完全驯服——即便手稿正确，它给出的也只是 Hausdorff 维数这一条战线上的胜利，极大函数、更高维、正测度之外的精细结构仍在原位。数学的进展向来如此：把一条防线推到理论上限，剩下的防线立刻变成新的前沿。这份稿件若被验证，它值得的位置是“四维战场的转折点”；若被修改或证伪，它同样已经证明了一件事——挂谷问题最前沿的攻击面，已经可以被机器系统性地测绘出来。

\begin{keybox}[可记忆主句]
对这份手稿，正确的姿态既不是膜拜也不是鄙视，而是把它放进验证流水线：先复演链条根部，再谈四维挂谷的前途；在证据到来之前，未知就是未知。
\end{keybox}